% =========================================================
% 2.2 RISK TERRAIN MODELING DIAGNOSTICS (RTMDX)
% El \section{} de esta sección está en el chapter.tex del capítulo;
% sus referencias van en seccion2_2.bib (misma carpeta).
% =========================================================

A diferencia del enfoque participativo y basado en percepciones de SafetiPin, el modelado del terreno de riesgo (RTM, por sus siglas en inglés) identifica los riesgos espaciales derivados de las características del entorno y modela cómo estos convergen para crear escenarios conductuales propicios para la delincuencia. El proceso de RTM comienza evaluando diversos factores que se consideran geográficamente relacionados con los incidentes delictivos\cite{caplan2016}.

Así, el software RTMDX no es más que una aplicación que automatiza la implementación del modelado RTM, la cual busca diagnosticar los factores espaciales de crimen en lugares con alto índice de criminalidad e identificar qué lugares son propicios para la generación de nuevos crímenes en el futuro. Este proceso se realiza mediante diez pasos, enumerados a continuación\cite{caplan2016}.

\begin{enumerate}
    \item Elegir el evento de resultado
    \item Elegir el área de estudio
    \item Elegir el periodo de tiempo
    \item Identificar los mejores factores de riesgo disponibles
    \item Obtener datos espaciales
    \item Mapear la influencia espacial de los factores
    \item Seleccionar los factores del modelo
    \item Ponderar los factores del modelo
    \item Combinar cuantitativamente los factores del modelo
    \item Comunicar información significativa
\end{enumerate}

Resulta relevante porque, hoy en día, permite auditar zonas utilizando únicamente la información disponible y, más importante aún, permite construir un modelo \enquote{predictivo}. Además, la naturaleza sistemática de sus diez pasos ofrece una guía metodológica clara que puede adaptarse al presente proyecto para estructurar la selección y ponderación de los factores de riesgo espacial considerados.
